\documentclass[14pt]{extarticle}
\input{external/preamble-latex-cools.tex}
\begin{document}
\begin{project}{Actividad de cierre}{Muffins en cajas.}{cool-muffinsEnCajas}%
Hay 30 \emph{muffins} y varias cajas para la feria de pastelería. En cada caja hay 6 \emph{muffins}. ¿Cuántas cajas hay?%
\begin{sidebyside}{2}{0.05}{0.05}{0.1}%
\begin{sbspanel}{0.4}%
Tyler escribió dos ecuaciones para este problema.%
\end{sbspanel}%
\begin{sbspanel}{0.4}%
\(\underline{\hspace{1 cm}} \times 6 = 30\)%
\par
\(30 \div 6 = \underline{\hspace{1 cm}}\)%
\end{sbspanel}%
\end{sidebyside}%
\par
Él dice que en cada espacio en blanco va el mismo número, aunque una ecuación es de multiplicación y la otra es de división. ¿Tiene razón? Explica o muestra tu razonamiento.%
\begin{image}{0}{1}{0}{}%
\includegraphics[max width=\linewidth, center]{external/whitespace-tikz/3cm.pdf}
\end{image}%
\end{project}
\end{document}